% ================================================================
\section{The Nyman--Beurling Distance as Vacuum Energy}
\label{sec:vacuum}
% ================================================================

\subsection{The Variational Principle}

The NB distance
\[
  d_N^2 = \inf_{v \in \RR^{N-1}} \int_0^1 (1 - f_{N,v})^2\,dx
  = 1 - b^T G^{-1} b
\]
where $b_k = \int_0^1 \fract{1/(kx)}\,dx = 1 - 1/k$ is the inner product
$\braket{h_k}{1}$, is obtained by minimizing over the weight vector $v$.

\begin{principle}[The Rayleigh--Ritz Principle]
This minimization is the \emph{Rayleigh--Ritz variational principle}:
we approximate the ground state energy by optimizing over a finite
trial space. The optimal weights $v_{\mathrm{opt}} = G^{-1}b$ are the
\emph{normal mode amplitudes} of the best approximation.

In the Cathedral, this is proved as \lean{nbDistSq\_formula}:
$d_N^2 = 1 - b^T G^{-1} b$. The physics: the vacuum energy is the
difference between the ``bare'' energy ($1$) and the energy extracted
by the optimal trial wavefunction ($b^T G^{-1} b$).
\end{principle}

The log-cutoff M\"obius witness $v_k = -\mu(k)(1 - \ln k / \ln N)$ used
in the forward direction is a \emph{Bartlett window}---a sub-optimal
trial wavefunction that extracts less energy than the exact minimizer.
In signal processing, this is the triangular apodization function.
In physics, it is a finite-temperature regularization of the
zero-temperature ground state.

\subsection{The Covariance Decomposition}

The Cathedral's Vasyunin tower decomposes the NB distance as:
\[
  d_N^2 = (1 - b^T v)^2 + v^T C v
\]
where $C = G - b b^T$ is the \emph{covariance matrix}.

\begin{correspondence}[Variance Decomposition]
This is the bias-variance decomposition in statistical learning:
\begin{itemize}
  \item $(1 - b^T v)^2$ is the \textbf{bias}---the squared distance
    from the mean. In physics: the \emph{mass term}.
  \item $v^T C v$ is the \textbf{variance}---the fluctuation energy.
    In physics: the \emph{kinetic energy} of quantum fluctuations.
\end{itemize}
RH is equivalent to both terms decaying to zero simultaneously:
the vacuum has zero bias \emph{and} zero fluctuation energy.
\end{correspondence}

